% from ega0-20-fr.tex
\oldpage[0\textsubscript{IV-1}]{118}
La prop. \sref{0.20.1.1-fr} s'interprète de la façon suivante :

\begin{corollary}[20.1.3]
\label{0.20.1.3-fr}
Étant donnés deux $A$-homomorphismes de $A$-anneaux $p:E\to C$,
$u:B\to C$ dont le premier a un noyau $\mathfrak{Q}$ de carré nul,
l'ensemble des $A$-homomorphismes $v:B\to E$ tels que $u=p\circ v$ est vide
ou est un espace homogène principal pour le groupe
$\operatorname{D\acute{e}r}_A(B,\mathfrak{Q})$.
\end{corollary}

En particulier :

\begin{corollary}[20.1.4]
\label{0.20.1.4-fr}
Soient $A$ un anneau, $C$ un $A$-anneau, $L$ un $C$-bimodule, $E$ une
$A$-extension de $C$ par $L$, $p:E\to C$ l'augmentation. L'application qui,
à toute dérivation $D\in\operatorname{D\acute{e}r}_A(C,L)$, associe l'application
$x\mapsto x+D(p(x))$ est un isomorphisme du groupe
$\operatorname{D\acute{e}r}_A(C,L)$ sur le groupe des $A$-équivalences de $E$ sur
elle-même.
\end{corollary}

\begin{proof}
Appliquons \sref{0.20.1.1-fr} pour $B=E$, $v_0=\mathrm{I}_E$ : l'ensemble
des $A$-homomorphismes $v:E\to E$ tels que $p\circ v=p$ est identique à
l'ensemble des applications $v=\mathrm{I}_E+D'$, où
$D'\in\operatorname{D\acute{e}r}_A(E,L)$. Dire qu'un tel $A$-homomorphisme $v$ est
une $A$-équivalence équivaut à dire que $v$ se réduit à l'injection canonique
dans $L$, ou encore que $D'(x)=0$ dans $L$ ; mais cela signifie aussi que
$D'$ se factorise en $E\xrightarrow{p}C\xrightarrow{D}L$, où $D$ est une
$A$-dérivation. D'où le corollaire.
\end{proof}

Pour les extensions triviales, \sref{0.20.1.3-fr} donne :

\begin{corollary}[20.1.5]
\label{0.20.1.5-fr}
Soient $A$ un anneau, $B$, $C$ deux $A$-anneaux, $L$ un $C$-bimodule,
$u:B\to C$ un $A$-homomorphisme ; l'application qui, à toute dérivation
$D\in\operatorname{D\acute{e}r}_A(B,L)$, associe l'application
$\alpha:x\mapsto(u(x),D(x))$ est une bijection sur l'ensemble des
$A$-homomorphismes $v:B\to D_C(L)$ (cf. 18.2.3) tels que $u$ se factorise en
$B\xrightarrow{v}D_C(L)\to C$.
\end{corollary}

Plus particulièrement :

\begin{corollary}[20.1.6]
\label{0.20.1.6-fr}
Soient $A$ un anneau, $B$ un $A$-anneau, $L$ un $B$-bimodule. Si, à toute
dérivation $D\in\operatorname{D\acute{e}r}_A(B,L)$, on associe : $1^\circ$ la
$A$-équivalence $(x,y)\mapsto(x,y+D(x))$ de l'extension $D_B(L)$ sur
elle-même ; $2^\circ$ le $A$-homomorphisme $x\mapsto(x,D(x))$ de $B$ dans
$D_B(L)$, inverse à droite de l'homomorphisme d'augmentation $D_B(L)\to B$,
on définit des correspondances biunivoques canoniques entre :
\begin{enumerate}
  \item[\rm(i)] l'ensemble $\operatorname{D\acute{e}r}_A(B,L)$ ;
  \item[\rm(ii)] l'ensemble des $A$-équivalences de $D_B(L)$ sur elle-même ;
  \item[\rm(iii)] l'ensemble des $A$-homomorphismes inverses à droite de
  l'homomorphisme d'augmentation $D_B(L)\to B$.
\end{enumerate}
\end{corollary}

On notera que la correspondance biunivoque ainsi établie entre (i) et (ii)
respecte les structures de groupe, et que celle qu'on en déduit entre (ii) et
(iii) n'est autre que la correspondance biunivoque déjà définie dans (18.3.8).

\begin{corollary}[20.1.7]
\label{0.20.1.7-fr}
Soient $A$ un anneau, $C$ un $A$-anneau, $L$ un $A$-bimodule, $E$ une
$A$-extension triviale de $C$ par $L$, $p:E\to C$ l'augmentation. L'ensemble
des $A$-dérivations $D$ de $E$ dans $L$ telles que $D(x)=x$ dans $L$ est
identique à l'ensemble des applications $x\mapsto x-w(p(x))$, où $w$
parcourt l'ensemble des $A$-homomorphismes inverses à droite de $p$.
\end{corollary}

\begin{proof}
Il suffit encore d'appliquer \sref{0.20.1.1-fr} pour $B=E$,
$v_0=\mathrm{I}_E$ ; si $v=\mathrm{I}_E-D$, la condition $D(x)=x$ pour
$x\in L$ équivaut à $v(x)=0$ pour $x\in L$, c'est-à-dire à $v=w\circ p$,
où $w:C\to E$ est un $A$-homomorphisme ; en outre la condition
$p\circ v=p$ équivaut à $p\circ w=\mathrm{I}_C$, autrement dit au fait que
$w$ est inverse à droite de $p$.
\end{proof}

\subsection{Propriétés fonctorielles des dérivations}
\label{subsection:0.20.2-fr}

